\section{Preliminaries}
\paragraph{Notation.} We will use big $O$ notations throughout the paper. Sometimes, it will be easier for us to use
$\ll$ and $\gg$ notations; when we say that $X\ll Y$ we mean that there is an absolute constant $C>0$ such that $X\leq C Y$.
Analogously, when we say that $X\gg Y$ we mean that there is an absolute constant $C>0$ such that $X\geq C Y$.

\subsection{Fourier Analysis over the Hypercube}
We consider the space of real-valued functions $f:\power{n} \to {\mathbb R}$,
equipped with the inner product $\inner{f}{g} = \Expect{x\in_R\power{n}}{f(x)g(x)}$.
It is well-known that the collection of functions $\set{\chi_S}_{S\subseteq[n]}$ defined by
$\chi_S(x) = \prod\limits_{i\in S}{(-1)^{x_i}}$ forms an orthonormal basis, so any function $f$
may be written as $f(x) = \sum\limits_{S\subseteq[n]}{\widehat{f}(S)\chi_S(x)}$ in a unique way,
where
$\widehat{f}(S) = \inner{f}{\chi_S}$.
We also consider $L_p$ norms for $p \geq 1$,
that are defined as $\norm{f}_p= \left(\Expect{x}{\card{f(x)}^p}\right)^{1/p}$.

For $z\in\power{n}$, we denote by $z_{-i}\in\power{n-1}$ the vector resulting from dropping
the $i$th coordinate of $z$. We denote by $(x_i = 1, x_{-i} = z_{-i})$ the $n$-bit vector
which is $1$ on the $i$th coordinate and is equal to $z$ on any other coordinate.
\begin{definition}
  The derivative of $f$ along direction $i\in[n]$
  is $\partial_i f\colon\power{n}\to\mathbb{R}$ defined by
  $\partial_i f(z) = f(x_i = 0, x_{-i} = z_{-i}) - f(x_i = 1, x_{-i} = z_{-i})$.
  The gradient of $f$, $\nabla f\colon\power{n}\to\mathbb{R}^n$, is defined
  as $\nabla f(z) = (\partial_1 f(z),\ldots,\partial_n f(z))$.
\end{definition}
Note that for a Boolean function $f$, for each $x\in\power{n}$, $\partial_i f(x) \neq 0$ if and only if the coordinate $i$
is sensitive on $x$, and in this case its value is either $1$ or $-1$. Thus, $\norm{\nabla f(x)}_2 = \sqrt{s_f(x)}$,
and it will be often convenient for us to use the $\ell_2$-norm of the gradient of $f$ in place of $\sqrt{s_f(x)}$.

\begin{fact}\label{fact:singeton}
  Let $f\colon\power{n}\to\mathbb{R}$, and let $i\in[n]$.
  Then $\partial_i f(x) =
  2\sum\limits_{S\ni i}{\widehat{f}(S)\chi_{S\setminus \set{i}}(x)}$.
  In particular, we have that $\Expect{x}{\partial_i f(x)} = 2\widehat{f}(\set{i})$.
\end{fact}


Next, we define the level $d$ \emph{Fourier weight} of a function $f$, the \emph{approximate level $d$ weight} of a function $f$
and the \emph{Fourier weight of $f$ up to/ above level $d$}.
\begin{definition}
  The level $d$ Fourier weight of $f$ is defined to be $W_{=d}[f] = \sum\limits_{\card{S} = d}{\widehat{f}(S)^2}$. The
  approximate level $d$ weight of a function $f$ is defined to be $W_{\approx d}[f] = \sum\limits_{d\leq j<2d}{W_{=j}[f]}$.
  The weight of $f$ up to level $d$ is defined as $W_{\leq d}[f] = \sum\limits_{\card{S}\leq d}{\widehat{f}(S)^2}$,
  and the weight of $f$ above level $d$ is defined as $W_{> d}[f] = \sum\limits_{\card{S}> d}{\widehat{f}(S)^2}$.
\end{definition}
For a function $f\colon\power{n}\to\mathbb{R}$, we also define the \emph{level at most $d$ part} of $f$, denoted by $f^{\leq d}$,
as $f^{\leq d}(x) = \sum\limits_{\card{S}\leq d} \widehat{f}(S)\chi_S(x)$. We note that by Parseval's equality we have that
$\norm{f^{\leq d}}_2^2 = W_{\leq d}[f]$.
\subsection{Random Restrictions}
For a function $f\colon\{0,1\}^n\to\mathbb{R}$, a set of coordinates $J\subseteq[n]$ and $z\in\power{\overline{J}}$, we define the
\emph{restriction} of
$f$ according to $(J,z)$ as the function $f_{\overline{J}\rightarrow z}\colon \{0,1\}^{J}\to\mathbb{R}$ defined by
\[
f_{\overline{J}\rightarrow z}(y) = f(x_{\overline{J}} = z, x_{J} = y).
\]
We refer to the coordinates of $J$ as ``alive'', and refer to the coordinates of $\overline{J}$ as fixed.
We will often be interested in random restrictions of a function $f$. By that, we mean that the set $J$ is chosen randomly by including in it each
$i\in[n]$ independently with some probability (which is specified), and $z\in\power{\overline{J}}$ is chosen uniformly.
We have the following fact, which establishes a relation between
the approximate level $d$ weight of $f$ and the level $1$ weight of a random restriction of $f$ in which each $i\in [n]$ is included in $J$ with probability
$\frac{1}{2d}$ (see~\cite[Fact 4.1]{kindler2018gaussian} and~\cite[Section 3.1]{KKKMS}).
\begin{fact}\label{fact:restriction}
  Let $f\colon\power{n}\to\mathbb{R}$ and $d\in\mathbb{N}$, let $(J,z)$ be a random restriction
  where each $j\in [n]$ is included in $J$ with probability $\frac{1}{2d}$, and $z\in\power{\bar{J}}$ is chosen uniformly. Then $\Expect{J,z}{W_{=1}[f_{\bar{J}\rightarrow z}]}\gg W_{\approx d}[f]$.
\end{fact}
\begin{proof}
  For a fixed $J$ and $T\subseteq J$, we have that
  $\widehat{f_{\bar{J}\rightarrow z}}(T) = \sum\limits_{S\subseteq \bar{J}}\widehat{f}(T\cup S)\chi_S(z)$. Thus,
  \[
  \Expect{J,z}{W_{=1}[f_{\bar{J}\rightarrow z}]}
  =\Expect{J}{\Expect{z}{\sum\limits_{j\in J}\left(\sum\limits_{S\subseteq\bar{J}}\widehat{f}(S\cup \set{j})\chi_S(z)\right)^2}}
  =\Expect{J}{\sum\limits_{j\in J}\sum\limits_{S\subseteq\bar{J}}\widehat{f}(S\cup \set{j})^2},
  \]
  where the last transition follows by pushing the expectation over $z$ inside and using Parseval's equality. It follows that
  \[
  \Expect{J,z}{W_{=1}[f_{\bar{J}\rightarrow z}]} =
  \Expect{J}{\sum\limits_{T}\widehat{f}(T)^21_{\card{T\cap J} = 1}}
  =\sum\limits_{T}\widehat{f}(T)^2\Prob{J}{\card{T\cap J}=1},
  \]
  and as $\Prob{J}{\card{T\cap J}=1}\gg 1$ for all $T$ whose size is between $d$ and $2d$ and the rest of the summands
  are all non-negative, the proof is concluded.
\end{proof}

\subsection{Bonami's Lemma}
The \emph{degree} of a function $f\colon \power{n}\to\mathbb{R}$ is defined as
${\sf deg}(f) = \max\sett{\card{S}}{\widehat{f}(S)\neq 0}$.
We will need the following consequence of
the Bonami-Beckner-Gross hypercontractive inequality~\cite{Bec75,Bonami,grossLogSob} (see also~\cite{ODonnell}), showing that the $4$-norm of a low-degree function $f$ is comparable to its $2$-norm.
\begin{thm}\label{thm:hypercontractivity}
  Let $f\colon\power{n}\to\mathbb{R}$ be a function of degree at most $d$. Then
  $\norm{f}_4\leq \sqrt{3}^d\norm{f}_2$.
\end{thm}
